% Introduction (front writer). Plan: ../OUTLINE.md "§1 Introduction"; synthesis research/notes-final.md §0, §1 (H1-H12, each
% item's "How to state it"), §2.3, §2.8, §5.3; research/final-check.md §4 (items 1, 2, 4, 5 bind the contributions list);
% the brief's agreed stances ("What a theory is", "Existing language": the ladder). Every claim carries the body's hypotheses
% or points to the result that states them. Credit: footnote F3 (OUTLINE §C) on the paragraph stating the two problems.
\section{Introduction}\label{sec:intro}

What does a scientific theory do, and what is it? On one view, the view behind Solomonoff's theory of inductive inference \citep{solomonoff1964formal}, a theory predicts observations and is a program: the ideal inductive reasoner is a Bayes mixture over all programs that predict the data, each with a weight exponentially small in its length. On another, a theory answers questions posed in some language, why-questions included \citep[ch.~5]{vanfraassen1980scientific}, and is a logical theory, a set of sentences closed under consequence from which the answers are derived \citep{suppe1977structure}. Other answers have been given to both questions: theories explain \citep{hempel1948studies} or solve problems \citep{laudan1977progress}; they are classes of models \citep{suppes1960comparison}, or practices of exemplars and skills \citep{kuhn1970structure,polanyi1958personal}.\footnote{Several of the questions studied here first appeared in H\"anni's working notes \citep{hanni2026notes}, among them whether theories are better seen as answering questions than as predicting raw data, and how a written theory relates to the reasoner who uses it.} This paper puts the second view into the form of the first. We study \SAI (\SAIabbr), a universal prior over axiom systems in the style of Solomonoff induction, and use it to make precise, and to test, claims about what theories do and what they are. The word ``axiom'' refers to the hypotheses, which are axiom systems, sets of first-order sentences produced by programs; it does not refer to the induction axiom of arithmetic.

\paragraph{The object.} Fix a first-order language $\mathcal L$ and a decidable consistent background theory $B$. The data are finitely many sentences labelled true or false. A universal machine reads a random tape, runs the program it finds at its start on the rest of the tape, and emits sentences, so the emitted theory is random. The prior probability $\MA(D)$ of a data set $D$ is the probability that this theory is consistent with $B$ and proves every datum: each sentence labelled true, and the negation of each sentence labelled false (\cref{def:setting:MA}). A hypothesis must explain each datum, not merely tolerate it. At a new sentence the posterior splits three ways: the random theory proves it, refutes it, or leaves it open. The open outcome is native, and with the probability $\Inc_0$ that the random theory is inconsistent it makes up the defect of the semimeasure $\MA$, which we keep rather than renormalise (\cref{rem:setting:defects}). The predictions form a belief function \citep{shafer1976mathematical}, updated by geometric conditioning; requiring only that a hypothesis be consistent with the data gives Dempster's rule instead, under which data eliminate hypotheses but never change the odds between the survivors (\cref{prop:basic:geometric,prop:basic:dempster,prop:basic:cr}).

\paragraph{Raw data and questions.} Solomonoff induction scores a hypothesis by how well it predicts a stream of observations, ideally raw. Science predicts much more particular things: where Mars is at a given time, to within two minutes of arc; the period of a planet, given the size of its orbit; whether an orbit is planar. These questions are posed with abstractions that are already available, such as positions, periods and orbits, which may themselves be highly theoretical. We order views of what a theory does on a ladder, each step further from predicting raw data: (1) predicting raw data, such as pixels or observation sequences; (2) answering questions posed with existing abstractions; (3) theories that extend the language with theoretical terms that pay for themselves, while the data stay in the base language; (4) theories that revise the base abstractions, as when a concept changes its meaning or a question is shown to be ill-posed. \SAIabbr over a fixed language sits on step 2. Its base language is whatever abstractions are available; it is not an observation language in the logical empiricists' sense, so the classical objections to a theory-neutral observation language do not apply to it in the same way, although theory-ladenness biases comparisons between theories that disagree about the base language (\cref{sec:phil:map}). We take step 2 to be a step in the right direction. We argue that it differs from step 1 in the target, the data type, logical coherence and partiality, not in the inference principle (\cref{sec:phil:questions,tab:phil:differences}): \SAIabbr is a Bayes mixture over a universal class, and version by version it agrees with Solomonoff operator induction restricted to consistent labellers (\cref{sec:soi}). We have partial results on step 3 (\cref{sec:shape,sec:terms}); step 4 has no formal handle here.

\paragraph{What a theory is.} On the second question the idealised formalism is neutral, and this is a theorem. Every consistent labeller of sentences, a program that answers some questions consistently, is an axiom system at a cost of $O(1)$ bits: coin-free \SAIabbr agrees with coin-free consistent \SOIabbr through Craig's construction inside the universal prior, and one-tape with one-tape, tape by tape (\cref{thm:soi:coinfree,thm:soi:onetape}); new vocabulary costs $O(1)$ bits (\cref{cor:basic:vocabulary,thm:terms:horn}); and program length sees no reasoner (\cref{thm:shape:invisible}). So \SAIabbr cannot see the difference between a theory and a program. We therefore keep two claims apart. That textbook presentations of fundamental theories are declarative, stating laws, definitions and models given by equations rather than procedures that print predictions, is an observation, with two qualifications: declarative is not sentential, and theories in use have procedural parts (\cref{sec:phil:observation}). Why that shape is good is an open question. Craig's construction rules out an answer based only on what can be represented, so any answer must be relative to a reasoner with costs, or to a community of reasoners (\cref{q:phil:declarative}); the best formal answer we have locates the large effects in the pair (reasoner, vocabulary), not in logical form (\cref{sec:shape}).

\paragraph{A preview: Kepler's third law.} Our running example is Kepler's laws of planetary motion. The third law, $\mathrm{K3}$, says that the squares of the periods of the planets are proportional to the cubes of the semi-major axes of their orbits. What it says about observations depends on the data language. In a language of heliocentric directions it adds nothing to Kepler's first two laws, since such a direction does not depend on the size of the orbit (\cref{prop:kepler:helio}). In a geocentric language, a modern idealisation of the data Kepler worked from, it has content: given the other elements of the orbits of the Earth and Mars, it fixes the size of the orbit of Mars and with it every geocentric direction of Mars, which without it ranges over an arc (\cref{prop:kepler:geo,prop:kepler:K3worth}). A computation puts what $\mathrm{K3}$ thus saves on the elements of Mars at $11.3$ to $14.7$ bits for precisions from $8'$ to $1'$, against a law text of $115$ to $375$ bits in the three codes we tried; so in none of these codes does compression alone select $\mathrm{K3}$ at Kepler's precision (an extrapolation from Mars to the other planets, conditional on the code; \cref{ex:kepler:K3cost}). And Kepler's laws are false: against the plan94 ephemeris a fixed ellipse misses Mars by $1.6'$ to $1.9'$ over $12$ to $30$ years (\cref{sec:kepler:wrong}). How a proof-requiring inducer should treat such a law is the subject of \cref{sec:laws}.

\subsection{Contributions}\label{sec:intro:contrib}

The main results, roughly in order of importance; the sections state them with their hypotheses, and \cref{sec:related} says what is new and what is not.
\begin{enumerate}[leftmargin=*,label=(\arabic*)]
\item \emph{\SAIabbr is consistent \SOIabbr, version by version:} coin-free with coin-free and one-tape with one-tape, within constant factors on every data set, and Kolmogorov with Kolmogorov within an additive constant; on labelled atoms of a fresh predicate \SAIabbr is \SOIabbr on strings, and the companion paper's results on binary labels that hold up to additive constants transfer (\cref{thm:soi:coinfree,thm:soi:onetape,thm:soi:transfer,cor:soi:transfer}).
\item \emph{Random theories that no r.e.\ theory can replace:} a stream on which the one-tape loss is bounded while no consistent r.e.\ theory proves the data and the limit posterior lives on random theories (\cref{thm:logic:lottery}); over a decidable true background containing Robinson arithmetic $\Q$, the data can all be true while the limit posterior is concentrated on false theories, a possibility result with adversarial data (\cref{cor:logic:false}); on $\Pi_1$ data over such a background the mixture is, up to a constant factor, one random true theory (\cref{thm:logic:threshold}).
\item \emph{Theory shape is relative to a reasoner and a vocabulary:} for Boolean laws, program length sees neither, while the least written length of a usable theory depends exponentially on the pair (reasoner, vocabulary) (\cref{thm:shape:invisible,cor:shape:regimes}); procedural form costs a factor $m$, the number of query directions of one parity law, over the base variables or with one round of reasoning, and at most a factor $9$ in written length with intermediate concepts and unbounded rounds (\cref{thm:shape:directions}).
\item \emph{The theoretician's dilemma inside an ideal inducer:} with free deduction, theoretical vocabulary changes observational predictions only by a constant factor; on literal data with charged deduction, padded (Craig-type) observational axiomatisations with polynomial-time membership fit exactly the polynomial-time regularities, vocabulary is eliminable at polynomial cost, and to fit more than the regular sets a finite theory needs a polyadic theoretical relation or objects unlike every named one (\cref{thm:terms:dilemma}).
\item \emph{Approximate laws:} over deterministic theories, proof-requiring compatibility rewards a false exact law over its true tolerance form; on data from a law with independent errors, the best loss rate is that of the law together with its error distribution, attained by $\MA$ when the distribution is computable (\cref{prop:laws:strength,prop:laws:errormodel}); on Kepler's residuals a law with a fitted error histogram wins every cell computed among the Kepler-based machines compared.
\item \emph{The basic theory of $\MA$:} dominance; a loss at most the complexity up to deductive equivalence, plus a constant, of any consistent r.e.\ theory that proves the data (semimeasure convention); an inconsistency probability $\Inc_0$ that is Solovay complete and random, so that $\MA$ is not lower semicomputable; no coding theorem (\cref{thm:basic:dominance,thm:basic:loss,thm:basic:inc,thm:basic:nocoding}).
\item \emph{Where \SAIabbr and \SOIabbr separate:} logical closure and an observed ``independent'' label separate them (\cref{prop:soi:closure,prop:soi:independent}); the separations under time bounds come from free deduction, not from axioms (\cref{thm:soi:diagonal,prop:soi:freededuction}); with deduction charged, \SAIabbr is certificate \SOIabbr (\cref{thm:soi:sandwich}).
\item \emph{G\"odel phenomena:} over a background containing $\Q$, belief or independent mass on a $\Pi_1$ sentence certifies its truth; over one containing $\PA$, the coin-free posterior never learns $\Sigma_1$-soundness; and the $\omega$-gap is the $\omega$-incompleteness of the background plus the instances (\cref{prop:logic:shape,thm:logic:soundness,prop:logic:omega}).
\item \emph{What a theory does:} answers cost no more description than raw modelling plus an extractor; in a toy world a conservation law's gain per raw step equals the raw no-law predictor's loss on the question about the conserved quantity, while its share of the raw loss vanishes (\cref{prop:phil:extractor,prop:phil:conservation}); deductively equivalent formulations keep their prior odds, and explanatory direction is recovered from simplicity only under a code that lists theories as written (\cref{prop:phil:flagpole,prop:phil:flagpolesystem}).
\item \emph{Best systems and explanatory power:} Lewis's ``kind nature'' is a kink in the hull of length against fit, and explanatory power is a per-bit log-likelihood ratio whose growth with the number of systems singles out laws (\cref{prop:laws:kink,def:laws:EP}).
\item \emph{Kepler and Newton:} over five historical forms for the orbit of Mars, on noise-free data the data decide compatibility and the syntax of the law texts decides among compatible forms, while with graded fit the data decide every step but one within tens of observations (\cref{tab:kepler:map,ex:kepler:graded}).
\end{enumerate}

\subsection{What we do not claim}\label{sec:intro:notclaim}

We do not claim that \SAIabbr shows a theory to be a logical theory rather than a program: in the idealised setting it cannot see the difference (\cref{sec:phil:craig}). That textbook theories are declarative is an observation, not a theorem, and why that shape is good is open (\cref{q:phil:declarative}). The results of \cref{sec:shape} hold in a stated model, with Boolean laws, named reasoners and written length as the cost; they are not theorems about textbooks, and whether they extend beyond Boolean laws is open (\cref{q:shape:beyond}). We say nothing realist about the theories an inducer favours: its limit posterior can live on random theories and, with true but adversarially chosen data, on false ones (\cref{thm:logic:lottery,cor:logic:false}). The proof-requiring core of \SAIabbr is the noise-free special case of the confrontation of theory with data; graded fit enters through randomised theories, not natively (\cref{sec:limits:stat}). \SAIabbr is an ideal of the evaluation of theories, not of their discovery, and it is not computable. Finally, our philosophical conclusions are arguments; \cref{sec:phil:map} says which of them rest on theorems.

\paragraph{Organisation.} \Cref{sec:setting,sec:basic} set up the definitions and the basic theory, \cref{sec:soi} relates \SAIabbr to operator induction, and \cref{sec:logic} studies what is specific to logic. \Cref{sec:phil} turns to what a theory does and is; \cref{sec:shape,sec:terms} treat theory shape and theoretical terms, \cref{sec:laws} approximate laws and best systems, and \cref{sec:kepler} the Kepler example. \Cref{sec:limits,sec:related,sec:discussion} discuss limitations, related work and open problems. Proofs and computations are in the appendices; \cref{app:notation} collects the notation.
